\noindent\hspace*{2em}神经网络计算图的多核执行能够提高并行度，但任务切分会同时改变跨域搬运、片上容量占用和共享带宽竞争。本文围绕独立任务、同核复用与共享只读缓存三种场景，建立从计算图切分到事件评价的统一调度方法，并在100张计算图、1至5核上完成可核验实验。

\textbf{针对问题一}，建立\textbf{Task驻留域结构搬运模型}。从生产者—消费者关系提取依赖边，构造整图、弱连通分量装箱和拓扑切块候选，以轻量评分安排评价顺序，再按Makespan优先选择最终方案。相对整图单核基准，五核平均加速比为\textbf{3.4383}；图001的工期由233110周期降至\textbf{47502}周期，新增搬运量为\textbf{4608字节}。

\textbf{针对问题二}，建立\textbf{核心驻留与物理副本容量模型}。将同核子图合并为核级Task，按物理副本从申请到最后读者完成的区间核算L1与UB占用，并把结构搬运、换出恢复和共享DDR排队纳入事件评价。五核平均加速比为\textbf{3.4738}，初始方案到最终方案的平均工期改善为\textbf{60.3622\%}；四核增至五核时，98张图的候选工期得到改善或保持。

\textbf{针对问题三}，建立\textbf{基于查询与完成事件的FIFO Cache模型}。按照发射查询、读取完成后插入和先进先出淘汰的时序，将Cache读与DDR读置于独立带宽服务中，并在问题二计划上比较硬件效应和调度适配。五核平均加速比为\textbf{3.5533}。在500组同核计划配对中，五核硬件、调度适配和综合效应比分别为\textbf{1.0146}、\textbf{1.0139}和\textbf{1.0289}。

固定问题二最终计划进行Cache参数敏感性实验：六张代表图的24组固定计划全部通过评价，37条重优化候选全部通过；容量减半、容量加倍、带宽减半和带宽加倍时，重优化工期相对默认配置的均值比分别为\textbf{0.9900}、\textbf{0.9869}、\textbf{0.9892}和\textbf{0.9885}。结果表明，Cache参数变化的影响取决于容量占用、查询时序与关键路径位置的共同作用。

\noindent{\zihao{3}\textbf{关键词：}}多核调度；张量驻留域；拓扑切块；容量约束；FIFO缓存
